\documentclass[12pt,oneside]{article}

% ---------------------------
% PAQUETES
% ---------------------------
\usepackage{geometry}
\usepackage[utf8]{inputenc}
\usepackage[spanish,es-nodecimaldot,es-tabla]{babel}
\usepackage{amsmath,amssymb}
\usepackage{microtype}
\usepackage[svgnames]{xcolor}
\usepackage{enumitem}
\usepackage{caption}
\usepackage{subcaption}
\usepackage{fancyhdr}
\usepackage{graphicx}
\usepackage{float}
\usepackage{listings}
\usepackage{mathptmx}
\usepackage{makecell}
\usepackage{multirow}
\usepackage{booktabs}
\usepackage{array}
\usepackage{longtable}
\usepackage{xcolor}
\usepackage{setspace}
\usepackage{colortbl}
\usepackage{pifont}

% ---------------------------
% MÁRGENES
% ---------------------------
\geometry{
  letterpaper,
  left=20mm,
  right=20mm,
  tmargin=30mm,
  bmargin=30mm
}

% ---------------------------
% COLORES
% ---------------------------
\definecolor{Custom_color_2}{rgb}{0.2,0.2,0.8}
\definecolor{RowLight}{RGB}{245,245,245}
\definecolor{CalifColor}{RGB}{122,74,46}
\definecolor{AsistColor}{RGB}{46,107,122}
\definecolor{AguaColor}{RGB}{30,100,180}
\definecolor{FisicaColor}{RGB}{59,130,80}
\definecolor{FuncColor}{RGB}{100,40,120}
\definecolor{CodeBg}{RGB}{245,245,245}
\definecolor{CodeFrame}{RGB}{180,180,180}

\captionsetup{justification=raggedright,labelfont={color=Custom_color_2,bf}}
\setlength{\headheight}{18pt}
\setlength{\parskip}{4mm}
\setlength{\parindent}{0mm}

% ---------------------------
% ENCABEZADO
% ---------------------------
\pagestyle{fancy}
\fancyhf{}
\renewcommand{\headrulewidth}{0pt}
\rfoot{\thepage}

% ---------------------------
% ESTILO DE CÓDIGO
% ---------------------------
\lstset{
  language=Python,
  basicstyle=\ttfamily\small,
  backgroundcolor=\color{CodeBg},
  frame=single,
  rulecolor=\color{CodeFrame},
  keywordstyle=\color{blue}\bfseries,
  stringstyle=\color{orange!80!black},
  commentstyle=\color{gray}\itshape,
  showstringspaces=false,
  breaklines=true,
  inputencoding=utf8,
  extendedchars=true,
  literate={á}{{\'a}}1 {é}{{\'e}}1 {í}{{\'i}}1 {ó}{{\'o}}1 {ú}{{\'u}}1
           {Á}{{\'A}}1 {É}{{\'E}}1 {Í}{{\'I}}1 {Ó}{{\'O}}1 {Ú}{{\'U}}1
           {ñ}{{\~n}}1 {Ñ}{{\~N}}1,
  numbers=left,
  numberstyle=\tiny\color{gray},
  numbersep=6pt,
  captionpos=b
}

% ---------------------------
% COMANDOS AUXILIARES
% ---------------------------
\newcommand{\appheader}[2]{%
  {\large\textbf{\textcolor{#1}{#2}}}%
  \vspace{1mm}\hrule height 1.2pt \vspace{3mm}%
}

\newcommand{\seccion}[1]{%
  \vspace{2mm}%
  \noindent\textbf{#1}%
  \vspace{1mm}%
}

\newcommand{\imagenresultado}[1]{%
  \begin{figure}[H]
    \centering
    \fbox{\begin{minipage}{0.92\textwidth}
      \vspace{4cm}
      \centering
      \textcolor{gray}{\textit{[Captura de pantalla -- #1]}}
      \vspace{4cm}
    \end{minipage}}
    \caption{Salida del programa -- #1}
  \end{figure}
}

% =========================================================
\begin{document}

% =========================================================
% PORTADA
% =========================================================
\Huge
\begin{center}
\begin{tabular}{ll}
\raisebox{-.4\height}{\includegraphics[scale=0.13]{Figures/Logo_ipn.png}} &
\textbf{ Instituto Politécnico Nacional}\\
\end{tabular}

\huge
\textbf{Unidad Profesional Interdisciplinaria de \\ Ingeniería Alejo Peralta}

\vspace{0.3cm}
\textbf{Licenciatura en Ciencia De Datos}

\textbf{Cómputo de Alto Desempeño}

\rule{\textwidth}{0.1pt}

\textbf{Actividad: Extracción, transformación y limpieza de datos}

\rule{\textwidth}{0.1pt}

\vspace{0.1cm}
\Large
\textbf{Dra. Leticia Oyuki Rojas Pérez} \\
\vspace{0.2cm}

\begin{center}
\begin{tabular}{l r}
\textbf{Integrante} & \textbf{Matrícula} \\[0.3cm]
Benjamín Gutiérrez Rivera & 2025720091 \\
\end{tabular}
\end{center}

\vspace{0.2cm}
\textbf{\today}
\end{center}

\newpage
\normalsize

% =========================================================
\tableofcontents
\newpage

% =========================================================
%  SECCIÓN: FUNCIONES DE CONVERSIÓN Y LIMPIEZA
% =========================================================
\section{Funciones de conversión y limpieza}

\appheader{FuncColor}{Descripción general}

Los cuatro programas usan funciones para convertir los datos del archivo a números.
Cada dato pasa por varios intentos hasta obtener un valor válido o descartarse.

\seccion{1.1 Función base \texttt{convertir} (todas las prácticas)}

Primero intenta convertir el texto a número con \texttt{float()}. Si no funciona
usa \texttt{text2num} para leer números escritos en español como \texttt{"ocho punto cinco"}.
Si ninguno de los dos funciona el dato se descarta.

\begin{lstlisting}[caption={Función base convertir}]
from text_to_num import text2num

def convertir(texto):
    try:
        return float(texto)           # intento 1: numero directo
    except ValueError:
        pass
    try:
        partes = texto.split(" punto ")
        numeros = [str(text2num(p.strip(), "es")) for p in partes]
        return float(".".join(numeros))   # intento 2: texto en espanol
    except (ValueError, KeyError):
        return None                   # dato invalido
\end{lstlisting}

\seccion{1.2 Diccionario y normalización de texto (Práctica 2)}

En la práctica 2 los datos pueden venir como palabras del tipo \texttt{Asistió} o \texttt{faltó}.
Se usa un diccionario para convertir esas palabras a \texttt{1.0} o \texttt{0.0}.
La función \texttt{normalizar} quita tildes y pone el texto en minúsculas antes de buscar
en el diccionario para que variantes como \texttt{ASISTIO} o \texttt{asistió} sean reconocidas.

\begin{lstlisting}[caption={Diccionario y normalización -- Práctica 2}]
import unicodedata

estado_asistencia = {
    "asistio":    1.0,
    "falto":      0.0,
    "no asistio": 0.0,
}

def normalizar(texto):
    texto = unicodedata.normalize("NFKD", texto.lower().strip())
    return "".join(c for c in texto if unicodedata.category(c) != "Mn")

def convertir(texto):
    clave = normalizar(texto)
    if clave in estado_asistencia:      # paso 0: buscar en diccionario
        return estado_asistencia[clave]
    try:
        return float(texto)
    except ValueError:
        pass
    try:
        partes = texto.split(" punto ")
        numeros = [str(text2num(p.strip(), "es")) for p in partes]
        return float(".".join(numeros))
    except (ValueError, KeyError):
        return None
\end{lstlisting}

\seccion{1.3 Limpieza de unidades -- litros (Práctica 3)}

Algunos datos vienen con la unidad escrita como \texttt{"1 litro"} o \texttt{"2 litros"}.
La función \texttt{limpiar} recorre una lista de unidades y las elimina del texto
antes de intentar la conversión.

\begin{lstlisting}[caption={Limpieza de unidades -- Práctica 3}]
unidades = ["litros", "litro"]

def limpiar(texto):
    for unidad in unidades:
        texto = texto.replace(unidad, "").strip()
    return texto

def convertir(texto):
    texto = limpiar(texto)
    try:
        return float(texto)
    except ValueError:
        pass
    try:
        partes = texto.split(" punto ")
        numeros = [str(text2num(p.strip(), "es")) for p in partes]
        return float(".".join(numeros))
    except (ValueError, KeyError):
        return None
\end{lstlisting}

\seccion{1.4 Limpieza de unidades -- minutos (Práctica 4)}

Igual que la práctica 3 pero con variantes de la unidad de tiempo como
\texttt{"min"}, \texttt{"mins"} y \texttt{"minutos"}.

\begin{lstlisting}[caption={Limpieza de unidades -- Práctica 4}]
unidades = ["minutos", "minuto", "mins", "min"]

def limpiar(texto):
    for unidad in unidades:
        texto = texto.replace(unidad, "").strip()
    return texto

def convertir(texto):
    texto = limpiar(texto)
    try:
        return float(texto)
    except ValueError:
        pass
    try:
        partes = texto.split(" punto ")
        numeros = [str(text2num(p.strip(), "es")) for p in partes]
        return float(".".join(numeros))
    except (ValueError, KeyError):
        return None
\end{lstlisting}

\seccion{1.5 Resumen del pipeline de conversión}

\begin{table}[H]
\centering
\caption{Pasos de conversión aplicados en cada práctica}
\renewcommand{\arraystretch}{1.4}
\begin{tabular}{
  >{\columncolor{FuncColor!10}}p{1.2cm}
  p{2.8cm}
  >{\columncolor{RowLight}}p{2.2cm}
  p{2.2cm}
  >{\columncolor{RowLight}}p{4.8cm}
}
\toprule
\rowcolor{FuncColor}
\textcolor{white}{\textbf{Paso}} &
\textcolor{white}{\textbf{Operación}} &
\textcolor{white}{\textbf{Práctica 1}} &
\textcolor{white}{\textbf{Práctica 2}} &
\textcolor{white}{\textbf{Prácticas 3 y 4}} \\
\midrule
0 & Diccionario       & ---           & \texttt{normalizar} + dict & --- \\
1 & Limpiar unidades  & ---           & ---        & \texttt{limpiar(texto)} \\
2 & \texttt{float()}  & \checkmark    & \checkmark & \checkmark \\
3 & \texttt{text2num} & \checkmark    & \checkmark & \checkmark \\
4 & Descartar         & \texttt{None} & \texttt{None} & \texttt{None} \\
\bottomrule
\end{tabular}
\end{table}

\seccion{1.6 Bibliotecas utilizadas}

\begin{table}[H]
\centering
\caption{Bibliotecas del proyecto}
\renewcommand{\arraystretch}{1.4}
\begin{tabular}{
  >{\columncolor{FuncColor!10}}p{3.5cm}
  p{4cm}
  >{\columncolor{RowLight}}p{5.8cm}
}
\toprule
\rowcolor{FuncColor}
\textcolor{white}{\textbf{Biblioteca}} &
\textcolor{white}{\textbf{Instalación}} &
\textcolor{white}{\textbf{Uso}} \\
\midrule
\texttt{pandas}        & \texttt{pip install pandas}   & Leer archivos y manejar los datos \\
\texttt{text\_to\_num} & \texttt{pip install text2num} & Convertir texto en español a número \\
\texttt{unicodedata}   & Estándar de Python            & Quitar tildes y normalizar texto \\
\bottomrule
\end{tabular}
\end{table}

\newpage

% =========================================================
%  PRÁCTICA 1: CALIFICACIONES
% =========================================================
\section{Práctica 1: Control de calificaciones}

\appheader{CalifColor}{Descripción general}

El programa lee el archivo \texttt{calificaciones.txt} donde cada línea tiene una calificación.
Los datos pueden ser números o palabras en español. Se usa la función \texttt{convertir}
de la sección~1.1 para procesarlos y luego se calculan estadísticas del grupo.

\seccion{2.1 Datos de entrada}

\begin{itemize}
  \item Archivo: \texttt{calificaciones.txt}
  \item Formato: un valor por línea como \texttt{8.5} u \texttt{ocho punto cinco}
  \item Columna generada: \texttt{calificacion} (número decimal)
\end{itemize}

\seccion{2.2 Cálculos realizados}

\begin{itemize}
  \item Total de registros válidos, promedio, máximo y mínimo
  \item Conteo de aprobados ($\geq 6.0$) y reprobados
  \item Nivel de desempeño del grupo con \texttt{if/elif/else}:
        Excelente ($\geq 9$), Bueno ($\geq 7$), Regular ($\geq 6$) o Requiere atención
\end{itemize}

\seccion{2.3 Tabla de métricas}

\begin{table}[H]
\centering
\caption{Métricas calculadas -- Práctica 1}
\renewcommand{\arraystretch}{1.4}
\begin{tabular}{
  >{\columncolor{CalifColor!15}}p{4.5cm}
  >{\columncolor{RowLight}}p{9cm}
}
\toprule
\rowcolor{CalifColor}
\textcolor{white}{\textbf{Métrica}} &
\textcolor{white}{\textbf{Descripción}} \\
\midrule
Total         & Calificaciones válidas leídas del archivo \\
Promedio      & Media de todas las calificaciones \\
Mayor / Menor & Calificación más alta y más baja \\
Aprobados     & Alumnos con calificación $\geq 6.0$ \\
Reprobados    & Alumnos con calificación $< 6.0$ \\
Desempeño     & Categoría del grupo según el promedio \\
\bottomrule
\end{tabular}
\end{table}

\seccion{2.4 Resultados obtenidos}

\begin{figure}[H]
  \centering
  \includegraphics[width=0.92\textwidth]{Figures/resultado_p1.png}
  \caption{Salida del programa -- Práctica 1: Control de calificaciones}
\end{figure}

\seccion{2.5 Conclusión}

El programa permite ver de forma rápida cómo está el grupo en general.
Con el promedio se puede saber si los alumnos están aprobando o no y qué tan lejos están
los mejores resultados de los más bajos. Esto ayuda a identificar si hay alumnos que
necesitan más apoyo.

\newpage

% =========================================================
%  PRÁCTICA 2: ASISTENCIAS
% =========================================================
\section{Práctica 2: Control de asistencia}

\appheader{AsistColor}{Descripción general}

El programa lee \texttt{asistencias.txt} donde cada línea indica si el alumno fue a clase o no.
Los valores pueden ser números como \texttt{1} o \texttt{0} o palabras como \texttt{Asistió}
y \texttt{faltó}. Se usa el diccionario de la sección~1.2 para manejar los casos de texto.

\seccion{3.1 Datos de entrada}

\begin{itemize}
  \item Archivo: \texttt{asistencias.txt}
  \item Formato: \texttt{1}/\texttt{0} o \texttt{Asistió}/\texttt{faltó}
  \item Columna generada: \texttt{asistencia} donde \texttt{1.0} es asistió y \texttt{0.0} no asistió
\end{itemize}

\seccion{3.2 Cálculos realizados}

\begin{itemize}
  \item Total de clases registradas
  \item Número de asistencias y faltas con \texttt{value\_counts()}
  \item Porcentaje de asistencia: $(asistencias / total) \times 100$
  \item Si el porcentaje es $\geq 80\%$ el alumno tiene derecho a evaluación
\end{itemize}

\seccion{3.3 Tabla de métricas}

\begin{table}[H]
\centering
\caption{Métricas calculadas -- Práctica 2}
\renewcommand{\arraystretch}{1.4}
\begin{tabular}{
  >{\columncolor{AsistColor!15}}p{4.5cm}
  >{\columncolor{RowLight}}p{9cm}
}
\toprule
\rowcolor{AsistColor}
\textcolor{white}{\textbf{Métrica}} &
\textcolor{white}{\textbf{Descripción}} \\
\midrule
Total de clases & Registros válidos leídos del archivo \\
Asistencias     & Días con valor \texttt{1.0} \\
Faltas          & Días con valor \texttt{0.0} \\
Porcentaje      & $(asistencias / total) \times 100$ \\
Estado          & Tiene o no derecho a evaluación \\
\bottomrule
\end{tabular}
\end{table}

\seccion{3.4 Resultados obtenidos}

\begin{figure}[H]
  \centering
  \includegraphics[width=0.92\textwidth]{Figures/resultado_p2.png}
  \caption{Salida del programa -- Práctica 2: Control de asistencia}
\end{figure}

\seccion{3.5 Conclusión}

Con este programa se puede saber si un alumno cumple el mínimo de asistencia para
presentar evaluación. Un porcentaje bajo puede significar que el alumno no ha
tenido suficiente contacto con el tema para ser evaluado de manera justa.

\newpage

% =========================================================
%  PRÁCTICA 3: CONSUMO DE AGUA
% =========================================================
\section{Práctica 3: Análisis de consumo de agua}

\appheader{AguaColor}{Descripción general}

El programa lee \texttt{consumo\_agua.txt} con el registro diario de litros de agua tomados.
Algunos datos incluyen la palabra \texttt{litro} o \texttt{litros} junto al número
por lo que se usa la función \texttt{limpiar} de la sección~1.3 para quitarla antes de convertir.

\seccion{4.1 Datos de entrada}

\begin{itemize}
  \item Archivo: \texttt{consumo\_agua.txt}
  \item Formato: \texttt{1.5}, \texttt{1 litro} o \texttt{dos punto cinco}
  \item Columna generada: \texttt{litros} (número decimal)
\end{itemize}

\seccion{4.2 Cálculos realizados}

\begin{itemize}
  \item Total de registros, suma total, promedio, máximo y mínimo
  \item Días en que se tomaron $\geq 2.0$ litros (meta diaria)
  \item Recomendación según el promedio: consumo adecuado o aumentar el consumo
\end{itemize}

\seccion{4.3 Tabla de métricas}

\begin{table}[H]
\centering
\caption{Métricas calculadas -- Práctica 3}
\renewcommand{\arraystretch}{1.4}
\begin{tabular}{
  >{\columncolor{AguaColor!15}}p{4.5cm}
  >{\columncolor{RowLight}}p{9cm}
}
\toprule
\rowcolor{AguaColor}
\textcolor{white}{\textbf{Métrica}} &
\textcolor{white}{\textbf{Descripción}} \\
\midrule
Registros      & Días con dato válido \\
Total          & Litros sumados en todo el periodo \\
Promedio       & Promedio diario de consumo \\
Mayor / Menor  & Día con más y menos consumo \\
Meta cumplida  & Días con $\geq 2.0$ litros \\
Recomendación  & Resultado según el promedio obtenido \\
\bottomrule
\end{tabular}
\end{table}

\seccion{4.4 Resultados obtenidos}

\begin{figure}[H]
  \centering
  \includegraphics[width=0.92\textwidth]{Figures/resultado_p3.png}
  \caption{Salida del programa -- Práctica 3: Análisis de consumo de agua}
\end{figure}

\seccion{4.5 Conclusión}

El programa muestra si una persona está tomando suficiente agua cada día.
Con el promedio y los días en que se cumplió la meta se puede ver qué tan
constante es el hábito y en qué días hace falta mejorar.

\newpage

% =========================================================
%  PRÁCTICA 4: ACTIVIDAD FÍSICA
% =========================================================
\section{Práctica 4: Registro de actividad física}

\appheader{FisicaColor}{Descripción general}

El programa lee \texttt{activida\_fisica.txt} con los minutos de ejercicio registrados por día.
Algunos datos incluyen la unidad escrita como \texttt{"60 min"} o \texttt{"30 minutos"}
por lo que se usa la función \texttt{limpiar} de la sección~1.4 para quitarla antes de convertir.

\seccion{5.1 Datos de entrada}

\begin{itemize}
  \item Archivo: \texttt{activida\_fisica.txt}
  \item Formato: \texttt{30}, \texttt{60 min} o \texttt{treinta minutos}
  \item Columna generada: \texttt{minutos} (número decimal)
\end{itemize}

\seccion{5.2 Cálculos realizados}

\begin{itemize}
  \item Total de registros, suma total, promedio, máximo y mínimo
  \item Días sin actividad (valor igual a $0$)
  \item Nivel de actividad con \texttt{if/elif/else}:
        Adecuada ($\geq 30$ min), Moderada ($\geq 15$ min) o Baja
\end{itemize}

\seccion{5.3 Tabla de métricas}

\begin{table}[H]
\centering
\caption{Métricas calculadas -- Práctica 4}
\renewcommand{\arraystretch}{1.4}
\begin{tabular}{
  >{\columncolor{FisicaColor!15}}p{4.5cm}
  >{\columncolor{RowLight}}p{9cm}
}
\toprule
\rowcolor{FisicaColor}
\textcolor{white}{\textbf{Métrica}} &
\textcolor{white}{\textbf{Descripción}} \\
\midrule
Registros          & Días con dato válido \\
Total              & Minutos sumados en todo el periodo \\
Promedio           & Promedio diario de minutos de ejercicio \\
Mayor / Menor      & Día con más y menos actividad \\
Días sin ejercicio & Días con valor igual a $0$ minutos \\
Nivel              & Adecuada, Moderada o Baja según el promedio \\
\bottomrule
\end{tabular}
\end{table}

\seccion{5.4 Resultados obtenidos}

\begin{figure}[H]
  \centering
  \includegraphics[width=0.92\textwidth]{Figures/resultado_p4.png}
  \caption{Salida del programa -- Práctica 4: Registro de actividad física}
\end{figure}

\seccion{5.5 Conclusión}

El programa permite ver si una persona hace ejercicio de forma regular.
Los días sin actividad y el promedio diario indican qué tan constante es el hábito.
Si el nivel es bajo o moderado conviene establecer una rutina más seguida para
llegar a los 30 minutos diarios recomendados.

\end{document}
